\documentclass[12pt, a4paper]{report}

%  Pacchetti 
\usepackage[T1]{fontenc}
\usepackage[utf8]{inputenc}
\usepackage[italian]{babel}
\usepackage{lmodern}
\usepackage{geometry}
\usepackage{graphicx}
\usepackage{xcolor}
\usepackage{listings}
\usepackage{booktabs}
\usepackage{longtable}
\usepackage{array}
\usepackage{tabularx}
\usepackage{hyperref}
\usepackage{fancyhdr}
\usepackage{titlesec}
\usepackage{enumitem}
\usepackage{amsmath}
\usepackage{amssymb}
\usepackage{tcolorbox}
\usepackage{mdframed}
\usepackage{microtype}
\usepackage{caption}
\usepackage{float}
\usepackage{parskip}

%  Layout pagina 
\geometry{
  top=2.5cm, bottom=2.5cm,
  left=3cm,  right=2.5cm,
  headheight=15pt
}

%  Colori 
\definecolor{codeblue}{RGB}{0, 80, 160}
\definecolor{codegray}{RGB}{128, 128, 128}
\definecolor{codegreen}{RGB}{0, 120, 50}
\definecolor{codebg}{RGB}{248, 248, 248}
\definecolor{codeborder}{RGB}{210, 210, 210}
\definecolor{noteblue}{RGB}{230, 241, 251}
\definecolor{noteblueborder}{RGB}{24, 95, 165}
\definecolor{warnbg}{RGB}{250, 238, 218}
\definecolor{warnborder}{RGB}{133, 79, 11}
\definecolor{okbg}{RGB}{234, 243, 222}
\definecolor{okborder}{RGB}{59, 109, 17}
\definecolor{headcolor}{RGB}{30, 50, 90}

%  Stile codice 
\lstdefinestyle{python}{
  language=Python,
  backgroundcolor=\color{codebg},
  basicstyle=\ttfamily\footnotesize,
  keywordstyle=\color{codeblue}\bfseries,
  stringstyle=\color{codegreen},
  commentstyle=\color{codegray}\itshape,
  numberstyle=\tiny\color{codegray},
  breaklines=true,
  breakatwhitespace=true,
  frame=single,
  rulecolor=\color{codeborder},
  numbers=left,
  numbersep=8pt,
  showstringspaces=false,
  tabsize=4,
  captionpos=b,
  literate={à}{{\`a}}1 {è}{{\`e}}1 {é}{{\'e}}1 {ì}{{\`i}}1
           {ò}{{\`o}}1 {ù}{{\`u}}1 
}
\lstset{style=python}

%  Box informativi 
\tcbuselibrary{skins, breakable}

\newtcolorbox{notebox}{
  colback=noteblue, colframe=noteblueborder,
  fonttitle=\bfseries\small, title=Nota,
  left=6pt, right=6pt, top=4pt, bottom=4pt,
  breakable
}
\newtcolorbox{warnbox}{
  colback=warnbg, colframe=warnborder,
  fonttitle=\bfseries\small, title=Attenzione,
  left=6pt, right=6pt, top=4pt, bottom=4pt,
  breakable
}
\newtcolorbox{okbox}{
  colback=okbg, colframe=okborder,
  fonttitle=\bfseries\small,
  left=6pt, right=6pt, top=4pt, bottom=4pt,
  breakable
}

%  Header/Footer 
\pagestyle{fancy}
\fancyhf{}
\fancyhead[L]{\small\textcolor{headcolor}{\textbf{Progetto Rientr@}}}
\fancyhead[R]{\small\textcolor{codegray}{Documentazione tecnica -- rdb2rdf\_step2.py}}
\fancyfoot[C]{\small\thepage}
\renewcommand{\headrulewidth}{0.4pt}

%  Titoli sezioni 
\titleformat{\chapter}[display]
  {\normalfont\Large\bfseries\color{headcolor}}
  {\chaptertitlename\ \thechapter}{10pt}{\Huge}
\titleformat{\section}
  {\normalfont\large\bfseries\color{headcolor}}
  {\thesection}{1em}{}
\titleformat{\subsection}
  {\normalfont\normalsize\bfseries}
  {\thesubsection}{1em}{}

%  Hyperref 
\hypersetup{
  colorlinks=true,
  linkcolor=codeblue,
  urlcolor=codeblue,
  citecolor=codeblue,
  pdftitle={Documentazione tecnica rdb2rdf\_step2.py},
  pdfauthor={Progetto Rientr@},
}

% 
\begin{document}
% 

%  Frontespizio 
\begin{titlepage}
  \centering
  \vspace*{2cm}
  {\Large\textbf{Progetto Rientr@Returns}\\[0.4em]
   \large CNR-STIIMA}\par
  \vspace{2.5cm}
  {\Huge\bfseries\color{headcolor} Documentazione Tecnica\par}
  \vspace{0.8cm}
  {\LARGE\ttfamily rdb2rdf\_step2.py\par}
  \vspace{0.5cm}
  {\large Matching tra DB relazionale e ontologia\\
   tramite String Matching e modelli NLP\par}
  \vspace{3cm}
  \begin{tabular}{ll}
    \textbf{Versione}   & Step 2 \\[0.3em]
    \textbf{Linguaggio} & Python 3.10+ \\[0.3em]
    \textbf{Ontologia}  & Rientra.rdf (namespace JobList) \\[0.3em]
    \textbf{Database}   & PostgreSQL 17 \\[0.3em]
  \end{tabular}
  \vfill
  {\small\textcolor{codegray}{Anno accademico 2025--2026}}
\end{titlepage}

%  Indice 
\tableofcontents
\clearpage

% 
\chapter{Introduzione e contesto}
% 

\section{Scopo del modulo}

Il modulo \texttt{rdb2rdf\_step2.py} costituisce il secondo passo del
processo RDB2RDF del progetto \textbf{Rientr@Returns}, sviluppato
presso CNR-STIIMA. L'obiettivo è risolvere il \emph{mismatch} tra le
denominazioni dei lavori presenti in un database relazionale esterno
(\texttt{ext\_job}) e le professioni rappresentate nell'ontologia
principale Rientr@, costruita sul vocabolario O*NET.

Il problema è formalmente definibile come un task di \emph{entity
linking}: dato un'entità (job title + descrizione) proveniente da una
sorgente eterogenea, trovare la classe ontologica corrispondente
nell'ontologia di riferimento, tenendo conto che i nomi possono essere
completamente diversi pur descrivendo la stessa professione.

\begin{notebox}
Esempio concreto del mismatch: il job title \textit{``Manager of the
digital archive''} nel DB relazionale corrisponde alla classe O*NET
\textit{``Archivists''} nell'ontologia. I due nomi non condividono
nessuna parola, ma la descrizione delle attività è semanticamente
equivalente.
\end{notebox}

\section{Posizione nel progetto}

Il modulo si colloca tra due fasi:

\begin{itemize}
  \item \textbf{A monte:} la tabella \texttt{ext\_job} in PostgreSQL,
        creata con \texttt{rientra\_ext\_jobs.sql}, contenente 9
        professioni con titolo e descrizione.
  \item \textbf{A valle:} un file excel
        (\texttt{rientra\_matching.xlsx}), che mostra i lavori del db con i match trovati nell'ontologia; vengono mostrati anche i risultati prodotti dalle strategie implementante per trovare i match.
\end{itemize}

% 
\chapter{Architettura e struttura del codice}
% 

\section{Panoramica della pipeline}

Il modulo esegue le seguenti fasi in sequenza:

\begin{enumerate}
  \item \textbf{Setup:} connessione a PostgreSQL, caricamento
        dell'ontologia RDF, estrazione dinamica delle professioni O*NET.
  \item \textbf{Lettura DB:} recupero dei job dalla tabella
        \texttt{ext\_job}.
  \item \textbf{Strategia 1:} String Matching lessicale
        (Jaccard + Overlap + Levenshtein).
  \item \textbf{Strategia 2:} NLP con BERT base,
        embedding \texttt{[CLS]}, cosine similarity.
  \item \textbf{Strategia 3:} NLP con BioBERT,
        embedding \texttt{[CLS]}, cosine similarity.
  \item \textbf{Strategia 4:} NLP con MiniLM tramite
        \texttt{sentence-transformers}, cosine similarity.
  \item \textbf{Confronto:} tabella comparativa dei risultati
        con verdetto di accordo tra modelli.
  \item \textbf{Output Excel:} creazione del file excel per la visualizzazione dei match trovati.
\end{enumerate}

\section{Organizzazione del file}

Il file è organizzato in sezioni logiche separate da commenti
delimitatori. La tabella seguente riassume le funzioni principali con
il loro ruolo.

\begin{longtable}{p{5.5cm} p{2.5cm} p{6cm}}
\toprule
\textbf{Funzione} & \textbf{Righe approx.} & \textbf{Ruolo} \\
\midrule
\endhead
\texttt{score\_color()} & 87--93 & Mappa score $\to$ stile rich \\
\texttt{score\_bar()} & 94--96 & Genera barra testuale ASCII \\
\texttt{get\_connection()} & 103--104 & Apre connessione PostgreSQL \\
\texttt{fetch\_all()} & 106--110 & Esegue query, restituisce lista dict \\
\texttt{test\_connection()} & 112--122 & Verifica connettività DB \\
\texttt{extract\_ontology\_jobs()} & 129--160 & Estrae professioni dall'ontologia \\
\texttt{\_build\_onto\_text()} & 162--169 & Costruisce testo composito O*NET \\
\texttt{\_norm(), \_jaccard(),} & 176--201 & Metriche string matching \\
\texttt{\_overlap(), \_levenshtein()} & & \\
\texttt{\_combined()} & 200--201 & Score pesato S1 \\
\texttt{s1\_match\_one()} & 203--215 & Match singolo job con S1 \\
\texttt{run\_string\_matching()} & 217--227 & Esegue S1 su tutti i job \\
\texttt{cosine\_similarity\_matrix()} & 234--253 & Cosine similarity vettorializzata \\
\texttt{get\_bert\_embedding()} & 260--288 & Embedding [CLS] con BERT raw \\
\texttt{run\_bert\_matching()} & 291--349 & Matching con BERT base o BioBERT \\
\texttt{run\_minilm\_matching()} & 356--407 & Matching con MiniLM \\
\texttt{print\_ontology\_table()} & 414--425 & Tabella rich professioni ontologia \\
\texttt{print\_db\_table()} & 428--438 & Tabella rich job nel DB \\
\texttt{print\_full\_comparison()} & 456--559 & Tabella comparativa strategie \\
\texttt{print\_summary\_multi()} & 609--639 & Pannello riepilogativo \\
\texttt{set\_header()} & 679--686 & Genera i setting \\
\texttt{style\_cell()} & 689--694 & Stile per il file excel \\
\texttt{save\_excel()} & 697--935 & Salva il file excel \\
\texttt{main()} & 722--847 & Orchestrazione pipeline \\
\bottomrule
\caption{Funzioni principali del modulo}
\end{longtable}

\section{Dipendenze}

\subsection{Dipendenze obbligatorie}

\begin{lstlisting}[caption={Installazione dipendenze minime}]
pip install psycopg2-binary rdflib scikit-learn rich numpy
\end{lstlisting}

\subsection{Dipendenze opzionali per NLP completo}

\begin{lstlisting}[caption={Installazione dipendenze NLP}]
pip install sentence-transformers transformers torch
\end{lstlisting}

\begin{notebox}
Il modulo implementa \emph{graceful degradation}: se
\texttt{transformers} o \texttt{sentence-transformers} non sono
installati, le relative strategie vengono saltate silenziosamente e
l'output contiene solo le strategie disponibili. 
\end{notebox}

% 
\chapter{Configurazione}
% 

\section{Parametri globali}

Tutti i parametri modificabili sono raccolti nella sezione di
configurazione in cima al file, in modo da non richiedere modifiche
al codice funzionale per adattare il modulo a un ambiente diverso.

\begin{lstlisting}[caption={Sezione di configurazione}]
DB_CONFIG = {
    "host":     "localhost",
    "port":     5432,
    "dbname":   "rientra_db",
    "user":     "postgres",
    "password": "postgres",
}

ONTOLOGY_FILE = "Rientra.rdf"
OUTPUT_DIR    = "output"

JOBLIST = Namespace("http://www.stiima.cnr.it/JobList#")
RIENTRA = Namespace("https://www.stiima.cnr.it/rientra#")

S1_THRESHOLD     = 0.12
BERT_THRESHOLD   = 0.70
MINILM_THRESHOLD = 0.50

MODELS = {
    "S2_BERT":    "bert-base-uncased",
    "S3_BioBERT": "dmis-lab/biobert-base-cased-v1.2",
    "S4_MiniLM":  "all-MiniLM-L6-v2",
}
\end{lstlisting}

\section{Soglie di confidenza}

Le soglie controllano quando un match viene considerato valido. Un
match con score inferiore alla soglia viene salvato come
\texttt{score\_raw} nell'output visivo ma \textbf{non viene inserito
nel grafo RDF} come tripla di collegamento.

\begin{table}[H]
\centering
\begin{tabular}{llll}
\toprule
\textbf{Strategia} & \textbf{Parametro} & \textbf{Valore attuale} &
\textbf{Note} \\
\midrule
S1 String Match  & \texttt{S1\_THRESHOLD}     & 0.12 & Bassa: metriche già combinate \\
S2 BERT base     & \texttt{BERT\_THRESHOLD}   & 0.70 & Da alzare a 0.92 (anisotropy) \\
S3 BioBERT       & \texttt{BERT\_THRESHOLD}   & 0.70 & Da alzare a 0.92 (anisotropy) \\
S4 MiniLM        & \texttt{MINILM\_THRESHOLD} & 0.50 & Calibrata, affidabile \\
\bottomrule
\end{tabular}
\caption{Soglie di confidenza configurate}
\end{table}

\begin{warnbox}
Le soglie di BERT base e BioBERT sono attualmente fissate a 0.70, ma
i risultati sperimentali mostrano che con questi modelli tutti i 9 job
superano la soglia indipendentemente dalla qualità del match, per via
del fenomeno di \emph{anisotropy} (si veda la Sezione~\ref{sec:anisotropy}).
Il valore raccomandato dal referente è $\geq 0.90$.
\end{warnbox}

% 
\chapter{Sorgenti dati}
% 

\section{Database relazionale esterno}

La tabella \texttt{ext\_job} in PostgreSQL costituisce la sorgente
dei job da matchare. Ha struttura intenzionalmente minimale: tre
colonne che rappresentano l'identificatore, il titolo e la
descrizione del lavoro.

\begin{lstlisting}[language=SQL, caption={Schema della tabella ext\_job}]
CREATE TABLE ext_job (
    id          VARCHAR(6)  PRIMARY KEY,  -- ext01, ext02, ...
    title       TEXT        NOT NULL,
    description TEXT
);
\end{lstlisting}

I 9 job attualmente presenti spaziano su domini molto diversi:
contabilità, gestione contenuti web, data science, relazioni
pubbliche digitali, operazioni di segreteria e vendita al dettaglio.
Questa eterogeneità è intenzionale e serve a testare i limiti del
vocabolario ontologico corrente.

\section{Ontologia Rientra.rdf}

L'ontologia viene letta dal file \texttt{Rientra.rdf} in formato
RDF/XML. Contiene 12.250 triple che definiscono la T-Box
(classi, proprietà, relazioni) del modello Rientr@.

\subsection{Namespace rilevanti}

\begin{table}[H]
\centering
\begin{tabular}{ll}
\toprule
\textbf{Prefisso} & \textbf{URI} \\
\midrule
\texttt{JobList:} & \texttt{http://www.stiima.cnr.it/JobList\#} \\
\texttt{rientra:} & \texttt{https://www.stiima.cnr.it/rientra\#} \\
\texttt{O1:}      & \texttt{http://www.stiima.cnr.it/O1:} \\
\bottomrule
\end{tabular}
\caption{Namespace principali dell'ontologia}
\end{table}

\subsection{Professioni O*NET presenti}

L'ontologia contiene attualmente 13 professioni nel namespace
\texttt{JobList}, estratte dalla classificazione O*NET:

\begin{table}[H]
\centering
\begin{tabular}{ll}
\toprule
\textbf{URI locale} & \textbf{Label O*NET} \\
\midrule
\texttt{Billing\_cost\_and\_rate\_clerks} & Billing, Cost, and Rate Clerks \\
\texttt{Bookkeping\_accounting\_and\_auditing\_clerks} & Bookkeeping,Accounting,Auditing Clerks \\
\texttt{Cabinet\_makers\_and\_bench\_carpenters} & Cabinetmakers and Bench Carpenters \\
\texttt{Construction\_laborers} & Construction laborers \\
\texttt{File\_clerks} & File Clerks \\
\texttt{Gem\_and\_diamond\_workers} & Gem and Diamond Workers \\
\texttt{Insurance\_claims\_clerks} & Insurance Claims Clerks \\
\texttt{Insurance\_policy\_processing\_clerks} & Insurance Policy Processing Clerks \\
\texttt{Landscaping\_and\_Groundskeeping\_Workers} & Landscaping and Groundskeeping Workers \\
\texttt{Postal\_service\_clerks} & Postal Service Clerks \\
\texttt{Receptionist\_and\_information\_clerks} & Receptionists and Information Clerks \\
\texttt{Travel\_guides} & Travel guides \\
\texttt{Word\_Processors\_and\_Typists} & Word Processors and Typists \\
\bottomrule
\end{tabular}
\caption{Professioni O*NET presenti nell'ontologia}
\end{table}

\subsection{Estrazione dinamica}

Le professioni non sono hardcodate nello script, ma vengono estratte
dinamicamente dall'ontologia tramite la funzione
\texttt{extract\_ontology\_jobs()}. Per ogni classe nel namespace
\texttt{JobList} vengono raccolti tre predicati:

\begin{itemize}
  \item \texttt{rdfs:label} -- nome canonico della professione
  \item \texttt{O1:Net\_job\_titles} -- titoli alternativi riportati
        da O*NET (es.\ ``Accounting Clerk, Accounting Specialist, ...'')
  \item \texttt{O1:Net\_short\_description} -- descrizione sintetica
        dell'attività lavorativa
\end{itemize}

Il prefisso ``Sample of reported job titles:'' viene rimosso
automaticamente da \texttt{Net\_job\_titles} prima del matching,
poiché altrimenti frammenti di testo descrittivo finirebbero nel
confronto lessicale causando match spuri.

% 
\chapter{Strategie di matching}
% 

\section{Strategia 1 -- String Matching}

\subsection{Descrizione}

La Strategia 1 è una baseline lessicale che non usa modelli di
linguaggio. Confronta il \textbf{titolo del job nel DB} con il label
canonico e tutti i titoli alternativi O*NET di ogni professione
nell'ontologia, usando tre metriche combinate linearmente.

\begin{notebox}
La Strategia 1 usa \textbf{solo il titolo} del job, non la
descrizione. Questo la rende veloce ma insensibile al contenuto
semantico delle attività lavorative.
\end{notebox}

\subsection{Metriche}

\subsubsection{Similarità di Jaccard}

\begin{equation}
J(A, B) = \frac{|A \cap B|}{|A \cup B|}
\end{equation}

dove $A$ e $B$ sono gli insiemi di parole (dopo normalizzazione in
minuscolo) dei due titoli. Misura la proporzione di parole condivise
sul totale delle parole distinte. Insensibile all'ordine delle parole.

\subsubsection{Overlap Coefficient}

\begin{equation}
O(A, B) = \frac{|A \cap B|}{\min(|A|, |B|)}
\end{equation}

Divide per la cardinalità del set più piccolo invece che per
l'unione. È più generoso di Jaccard quando un titolo breve è
semanticamente contenuto in uno più lungo (es.\ ``Data entry''
contenuto in ``Data Entry Operator'').
\newpage
\subsubsection{Similarità di Levenshtein normalizzata}

\begin{equation}
L(s_1, s_2) = 1 - \frac{d_{edit}(s_1, s_2)}{\max(|s_1|, |s_2|)}
\end{equation}

dove $d_{edit}$ è la distanza di edit (numero minimo di inserimenti,
cancellazioni e sostituzioni di caratteri). È l'unica delle tre
metriche sensibile all'ordine dei caratteri -- utile per abbreviazioni
e varianti ortografiche.

\subsubsection{Score combinato}

\begin{equation}
\text{score}_{S1}(a, b) = 0.5 \cdot J(a, b) + 0.3 \cdot O(a, b) + 0.2 \cdot L(a, b)
\end{equation}

I pesi assegnano priorità a Jaccard come metrica principale, con
Overlap come supporto per asimmetrie di lunghezza e Levenshtein per
variazioni ortografiche.

\subsection{Implementazione}

\begin{lstlisting}[caption={Funzione di matching S1}]
def s1_match_one(title: str, onto_jobs: list[dict]) -> dict | None:
    best, bj, bv = 0.0, None, None
    for j in onto_jobs:
        cands = [j["label"]]
        if j["titles"]:
            cands += [t.strip() for t in j["titles"].split(",") if t.strip()]
        for c in cands:
            sc = _combined(title, c)
            if sc > best:
                best, bj, bv = sc, j, c
    if best >= S1_THRESHOLD:
        return {"job": bj, "score": best, "via": bv}
    return None
\end{lstlisting}

La chiave \texttt{"via"} nel dizionario restituito indica esattamente
quale label specifica (tra quelle alternative O*NET) ha prodotto lo
score massimo -- informazione utile per il debugging e per la
valutazione della qualità del match.

% 
\section{Cosine Similarity -- funzione condivisa}

Tutte e tre le strategie NLP usano la stessa funzione per calcolare
la similarità coseno, garantendo uniformità metodologica nei confronti.

\begin{equation}
\cos(\theta) = \frac{\mathbf{A} \cdot \mathbf{B}}{\|\mathbf{A}\| \cdot \|\mathbf{B}\|}
\end{equation}

dove $\mathbf{A}$ è il vettore embedding del job nel DB e $\mathbf{B}$
è il vettore embedding di una professione O*NET. Il risultato è
compreso tra $-1$ e $1$; per embedding di modelli linguistici il
range effettivo è $[0, 1]$.
\newpage
\begin{lstlisting}[caption={Implementazione vettorializzata della cosine similarity}]
def cosine_similarity_matrix(query_emb: np.ndarray,
                              corpus_embs: np.ndarray) -> np.ndarray:
    query_norm  = query_emb / (np.linalg.norm(query_emb) + 1e-9)
    corpus_norm = corpus_embs / (
        np.linalg.norm(corpus_embs, axis=1, keepdims=True) + 1e-9
    )
    return corpus_norm.dot(query_norm)
\end{lstlisting}

Il termine $+10^{-9}$ (epsilon di stabilità numerica) previene la
divisione per zero nel caso teorico in cui un vettore abbia norma
nulla. La forma vettorializzata calcola in un'unica operazione la
similarità tra il vettore query e tutti i 13 vettori O*NET, senza
loop Python.

% 
\section{Strategia 2 -- BERT base}

\subsection{Modello}

\begin{table}[H]
\centering
\begin{tabular}{ll}
\toprule
\textbf{Parametro} & \textbf{Valore} \\
\midrule
Identificatore HuggingFace & \texttt{bert-base-uncased} \\
Numero di parametri & 110 milioni \\
Architettura & Transformer encoder, 12 layer \\
Dimensione embedding & 768 \\
Pre-training & Wikipedia EN + BookCorpus \\
Task di pre-training & Masked LM + Next Sentence Prediction \\
\bottomrule
\end{tabular}
\caption{Caratteristiche di bert-base-uncased}
\end{table}

\subsection{Strategia di embedding}

BERT non produce nativamente un embedding di frase. Per ricavarlo si
utilizza il vettore del token speciale \texttt{[CLS]}, inserito
automaticamente all'inizio di ogni sequenza di input. Questo token è
stato progettato durante il pre-training per catturare una
rappresentazione aggregata dell'intera sequenza.

\begin{lstlisting}[caption={Estrazione embedding [CLS] da BERT}]
def get_bert_embedding(text, tokenizer, model, max_length=512):
    inputs = tokenizer(text, return_tensors="pt",
                       truncation=True, max_length=max_length,
                       padding=True)
    with torch.no_grad():
        outputs = model(**inputs)
    # shape: [batch=1, seq_len, 768]  ->  [768]
    cls_embedding = outputs.last_hidden_state[:, 0, :].squeeze().numpy()
    return cls_embedding
\end{lstlisting}

\texttt{torch.no\_grad()} disabilita il calcolo del gradiente poiché
il modello è usato solo per inferenza, riducendo l'uso di memoria.

\subsection{Limitazioni: anisotropy}
\label{sec:anisotropy}

BERT non è stato fine-tuned per la sentence similarity. I suoi
embedding tendono a occupare una zona ristretta dello spazio
vettoriale (fenomeno noto come \emph{anisotropy}): i vettori puntano
tutti in direzioni simili, rendendo la cosine similarity
sistematicamente alta indipendentemente dalla reale similarità
semantica. In pratica, con soglia 0.70, tutti i 9 job del dataset
ottengono un match -- il che non è necessariamente corretto.

% 
\section{Strategia 3 -- BioBERT}

\subsection{Modello}

\begin{table}[H]
\centering
\begin{tabular}{ll}
\toprule
\textbf{Parametro} & \textbf{Valore} \\
\midrule
Identificatore HuggingFace & \texttt{dmis-lab/biobert-base-cased-v1.2} \\
Numero di parametri & 110 milioni \\
Architettura & Transformer encoder, 12 layer (come BERT) \\
Dimensione embedding & 768 \\
Pre-training & BERT base + PubMed abstracts + PMC full text \\
Corpus aggiuntivo & 4.5 miliardi di parole di letteratura biomedica \\
\bottomrule
\end{tabular}
\caption{Caratteristiche di BioBERT}
\end{table}

\subsection{Ragione dell'inclusione}

BioBERT è incluso per \textbf{confronto metodologico}, non perché
sia il modello migliore per questo task. La sua presenza consente di
osservare empiricamente come la specializzazione di dominio influenzi
i risultati: un modello addestrato su letteratura medica produce
embedding diversi da BERT base quando processa testi di ambito
lavorativo-amministrativo.

\subsection{Risultati osservati}

I test mostrano che BioBERT mappa 8 job su 9 sulla stessa classe
(\textit{File Clerks}), indipendentemente dal contenuto. Questo
comportamento è attribuibile al disallineamento tra il dominio di
addestramento (letteratura biomedica) e il dominio del task (job
matching). La classe \textit{File Clerks} ha una descrizione O*NET
breve e generica che produce un vettore ``neutro'' vicino a molti
testi in uno spazio semantico non calibrato per questo dominio.
\\
\begin{warnbox}
    L'analisi delle caratteristiche architetturali di BioBERT, unite alle prestazioni deludenti rilevate in fase sperimentale, suggerisce l'inefficacia del suo impiego nel caso in esame. Contrariamente a quanto osservato con il modello BERT standard, la specificità del pre-addestramento di BioBERT su domini biomedici non ha prodotto i benefici attesi, rendendo superflua e non proficua qualsiasi ulteriore attività di fine-tuning.
\end{warnbox}


% 
\section{Strategia 4 -- MiniLM}

\subsection{Modello}

\begin{table}[H]
\centering
\begin{tabular}{ll}
\toprule
\textbf{Parametro} & \textbf{Valore} \\
\midrule
Identificatore HuggingFace & \texttt{sentence-transformers/all-MiniLM-L6-v2} \\
Numero di parametri & 22 milioni \\
Architettura & Transformer encoder, 6 layer (distillato da BERT) \\
Dimensione embedding & 384 \\
Fine-tuning & Sentence similarity su 1+ miliardo di coppie \\
Libreria & \texttt{sentence-transformers} \\
\bottomrule
\end{tabular}
\caption{Caratteristiche di MiniLM}
\end{table}

\subsection{Differenze rispetto a BERT raw}

MiniLM è stato addestrato con un obiettivo esplicito di sentence
similarity tramite architettura \emph{Siamese}: due reti con pesi
condivisi processano coppie di frasi e vengono ottimizzate per
produrre vettori vicini per frasi semanticamente simili e vettori
lontani per frasi semanticamente diverse. Il risultato è che i suoi
score di cosine similarity sono \textbf{calibrati}: uno score di 0.86
indica genuina similarità semantica, uno score di 0.30 indica
genuina dissimilarità.

\subsection{Implementazione}

\begin{lstlisting}[caption={Matching con MiniLM}]
model = SentenceTransformer("all-MiniLM-L6-v2")
onto_embs = model.encode(onto_texts, convert_to_numpy=True)

for i, row in enumerate(db_jobs):
    q_emb = model.encode(db_texts[i], convert_to_numpy=True)
    # Cosine similarity calcolata con la stessa funzione di S2/S3
    sims  = cosine_similarity_matrix(q_emb, onto_embs)
    idx   = int(np.argmax(sims))
    sc    = float(sims[idx])
    m = {"job": onto_jobs[idx], "score": sc} if sc >= threshold else None
\end{lstlisting}

Gli embedding O*NET vengono pre-calcolati una sola volta prima del
loop sui job del DB, ottimizzando le prestazioni. Nonostante MiniLM
fornisca internamente funzioni di cosine similarity, la funzione
\texttt{cosine\_similarity\_matrix()} viene usata esplicitamente per
uniformità con S2 e S3.

\newpage
\subsection{Input testuale}

Sia per i job del DB che per le professioni O*NET, il testo passato
al modello è una concatenazione di più campi:

\begin{itemize}
  \item \textbf{Job nel DB:} \texttt{title + ". " + description}
  \item \textbf{Professione O*NET:} \texttt{label + Net\_job\_titles + Net\_short\_description}
\end{itemize}

Più contesto semantico viene fornito, più precisi risultano gli
embedding prodotti dal modello.


% 
\chapter{Produzione dei risultati in formato Excel}
% 

\section{Scopo dell'output Excel}

A differenza della fase precedente, il modulo attuale non genera un grafo RDF A-Box, ma produce un report analitico in formato Microsoft Excel (\texttt{.xlsx}). Questa scelta è motivata dalla necessità di validare le soglie di confidenza e confrontare le diverse strategie di matching (lessicali e neurali) in modo leggibile e strutturato prima della definitiva persistenza nel Knowledge Graph.

Il file viene generato automaticamente nella cartella \texttt{output/} con il nome\\ \texttt{rientra\_matching.xlsx}.

\section{Struttura del Workbook}

Il documento Excel è organizzato in tre fogli di lavoro (sheet) distinti, ognuno con un ruolo specifico nella pipeline di valutazione dei risultati:

\begin{enumerate}
  \item \textbf{Matching Results}: Foglio principale che contiene i risultati di tutte le strategie di matching per ogni job del database.
  \item \textbf{Ontology Jobs}: Elenco delle professioni O*NET estratte dall'ontologia originale, utilizzato come riferimento.
  \item \textbf{Score Detail}: Matrice dettagliata che riporta gli score raw per ogni combinazione job-modello.
\end{enumerate}

\section{Dettaglio del foglio ``Matching Results''}

Il foglio principale raccoglie le informazioni provenienti dal database PostgreSQL e le arricchisce con i risultati dei modelli NLP.

\subsection{Metadati e colonne}

Per ogni riga (corrispondente a un job del DB), vengono riportati:
\begin{itemize}
    \item \textbf{Dati DB}: ID (\texttt{ext\_job}), Titolo e Descrizione.
    \item \textbf{IRI Individuo}: L'URI che identificherebbe l'entità nel grafo (\texttt{rientra:ExtJob\_*}).
    \item \textbf{S1 String Match}: Label trovata, punteggio e etichetta specifica utilizzata per il match.
    \item \textbf{Strategie NLP (S2, S3, S4)}: Professione identificata, score di \textit{cosine similarity} e URI della classe corrispondente.
\end{itemize}

\subsection{Logica del Verdetto}

L'ultima colonna del foglio fornisce un \textbf{Verdetto} sintetico basato sul consenso tra i modelli NLP.

\begin{table}[H]
\centering
\begin{tabular}{lp{9cm}}
\toprule
\textbf{Verdetto} & \textbf{Descrizione} \\
\midrule
Tutti concordano & Tutti i modelli NLP attivi hanno identificato lo stesso URI O*NET. \\
Accordo parziale & Esiste un consenso tra alcuni modelli o lo score è considerato affidabile. \\
Divergono & I modelli NLP puntano a professioni O*NET differenti. \\
Solo S1 & Solo lo string matching ha prodotto un risultato sopra soglia. \\
Nessun match & Nessuna strategia ha superato le soglie di confidenza configurate. \\
\bottomrule
\end{tabular}
\caption{Logica del verdetto nel foglio Excel}
\end{table}

\section{Stili e Formattazione Condizionale}

Per facilitare l'analisi visiva, il modulo applica stili specifici basati sulla qualità del punteggio:

\begin{itemize}
    \item \textbf{Colorazione per Strategia}: Ogni colonna di matching ha un colore distintivo (es. giallo per S1, verde per MiniLM).
    \item \textbf{Evidenziazione Score}: I punteggi superiori alle soglie ($\geq 0.70$ per BERT e $\geq 0.50$ per MiniLM) sono evidenziati in grassetto verde.
    \item \textbf{Sotto Soglia}: I risultati insufficienti riportano il valore \textit{raw} in rosso con la dicitura ``sotto soglia''.
\end{itemize}
\newpage
\section{Esempio di implementazione}

Di seguito è riportata la struttura della funzione Python che utilizza la libreria \texttt{openpyxl}:

\begin{lstlisting}[caption={Snippet della funzione \texttt{save\_excel}}]
def save_excel(db_jobs, s1_res, nlp_results, model_names, onto_jobs):
    wb = Workbook()
    ws1 = wb.active
    ws1.title = "Matching Results"
    # Configurazione header e stili
    for r_idx, row in enumerate(db_jobs, start=2):
        # Scrittura dati DB e risultati strategie
        # Calcolo del verdetto di consenso
    wb.save(os.path.join(OUTPUT_DIR, "rientra_matching.xlsx"))
\end{lstlisting}

% 
\chapter{Output e interpretazione dei risultati}
% 

\section{Tabella comparativa}

La sezione 6 dell'output mostra una tabella con una riga per ogni
job del DB e una coppia di colonne (match + score) per ogni strategia
disponibile. La colonna \textbf{Verdetto} riassume l'accordo tra i
modelli NLP:

\begin{table}[H]
\centering
\begin{tabular}{lll}
\toprule
\textbf{Verdetto} & \textbf{Significato} & \textbf{Affidabilità} \\
\midrule
\texttt{\checkmark\ tutti}  & Tutti i modelli NLP concordano & Alta \\
\texttt{\textasciitilde\ parz.}  & Accordo parziale (alcuni modelli) & Media \\
\texttt{$\neq$\ div.}   & I modelli puntano a professioni diverse & Bassa \\
\texttt{S1 only}  & Solo string matching ha un match & Verificare manualmente \\
\texttt{$\times$\ no}     & Nessun match sopra soglia & Job non presente in ontologia \\
\bottomrule
\end{tabular}
\caption{Interpretazione della colonna Verdetto}
\end{table}

\begin{notebox}
Il verdetto è calcolato confrontando gli URI delle professioni trovate,
non gli score. Due modelli con score molto diversi (es.\ 0.929 e 0.863)
che puntano alla stessa professione danno verdetto ``\texttt{\checkmark\ tutti}''.
Due modelli con score simili che puntano a professioni diverse danno
``\texttt{$\neq$\ div.}''.
\end{notebox}

\section{Risultati correnti}

Con le 13 professioni attualmente nell'ontologia e le soglie
configurate:

\begin{table}[H]
\centering
\begin{tabular}{llll}
\toprule
\textbf{Job} & \textbf{S4 MiniLM} & \textbf{Match} & \textbf{Affidabilità} \\
\midrule
ext01 Accounting clerk & 0.863 & Bookkeeping Clerks & Alta \\
ext05 Data entry employee & 0.556 & Bookkeeping Clerks & Discutibile \\
ext09 Supermarket cashier & 0.511 & Billing Clerks & Discutibile \\
ext02--04, 06--08 & $< 0.50$ & Nessun match & Professione assente \\
\bottomrule
\end{tabular}
\caption{Risultati MiniLM con soglia 0.50}
\end{table}

Il solo match semanticamente affidabile è \texttt{ext01}. Gli altri
6 job non trovano corrispondenza perché le loro professioni
(Data analyst, Data scientist, Content manager web, ecc.) non sono
presenti tra le 13 classi dell'ontologia.

% 
\chapter{Limitazioni note e sviluppi futuri}
% 

\section{Limitazioni correnti}

\subsection{Vocabolario ontologico insufficiente}

Il collo di bottiglia principale non è la qualità dei modelli NLP ma
la limitata copertura del vocabolario O*NET nell'ontologia: 13
professioni su migliaia disponibili. L'aggiunta delle professioni
mancanti (Data analyst, Data scientist, Content manager, ecc.)
migliorerà drasticamente i risultati di MiniLM.

\subsection{Soglie BERT non calibrate}

Con soglia 0.70, BERT e BioBERT trovano un match per tutti i 9 job
indipendentemente dalla qualità. La soglia va alzata a $\geq 0.90$
come raccomandato dal referente del progetto.

\subsection{BioBERT -- dominio errato}

BioBERT è stato incluso per confronto metodologico ma non è adatto
al task di job matching. Produce risultati non affidabili (8/9 job
mappati su \textit{File Clerks}) e potrebbe essere sostituito con un
modello più appropriato per il dominio lavorativo.

\section{Sviluppi pianificati}

\begin{enumerate}
  \item \textbf{Aggiornamento soglie:} portare
        \texttt{BERT\_THRESHOLD} a 0.92 e aggiungere la colonna
        ``gap'' (differenza tra primo e secondo score) come misura di
        robustezza del match.
  \item \textbf{Arricchimento ontologia:} aggiungere le 10 nuove
        professioni O*NET che coprono i job attualmente senza match.

  \item \textbf{Modello modificato:} Si è scelto di superare la mappatura statica tra i framework O*NET ed ESCO, optando per l'implementazione di un modello BERT addestrato su dataset reali. Tale approccio è finalizzato all'individuazione accurata di equivalenze e similitudini semantiche tra i profili.
\end{enumerate}

% 
\chapter{Guida all'uso}
% 

\section{Prerequisiti}

\begin{enumerate}
  \item PostgreSQL 17 attivo e accessibile su \texttt{localhost:5432}
  \item Database \texttt{rientra\_db} creato:
        \lstinline|createdb -U postgres rientra_db|
  \item Tabella \texttt{ext\_job} popolata:
        \lstinline|psql -U postgres -d rientra_db -f rientra_ext_jobs.sql|
  \item File \texttt{Rientra.rdf} nella stessa directory dello script
  \item Ambiente virtuale Python attivato
\end{enumerate}

\section{Installazione}

\begin{lstlisting}[caption={Setup ambiente}]
python3 -m venv venv
source venv/bin/activate          # macOS/Linux
# oppure: venv\Scripts\activate   # Windows

pip install psycopg2-binary rdflib scikit-learn rich numpy
pip install transformers torch sentence-transformers
\end{lstlisting}

\section{Esecuzione}

\begin{lstlisting}[caption={Avvio dello script}]
python3 rdb2rdf_step2.py # macOS

python rdb2rdf_step2.py # Windows
\end{lstlisting}

\section{Output atteso}

Al termine dell'esecuzione viene prodotto nella cartella
\texttt{output/}:

\begin{itemize}
  \item \texttt{rientra\_matching.xlsx} -- file excel 
\end{itemize}

L'esecuzione completa con tutti i modelli NLP richiede circa 2--5
minuti alla prima esecuzione (download dei modelli da HuggingFace) e
circa 30--60 secondi nelle esecuzioni successive (cache locale).

% 
\appendix
\chapter{Messaggi di warning noti}
% 

Durante l'esecuzione compaiono alcuni messaggi di warning che sono
normali e non indicano errori:

\begin{table}[H]
\centering
\begin{tabular}{p{6cm} p{8cm}}
\toprule
\textbf{Messaggio} & \textbf{Causa e comportamento} \\
\midrule
\texttt{cls.predictions.* | UNEXPECTED} &
Layer di pre-training (Masked LM) non usati. Vengono scartati senza
effetti sul funzionamento. \\[0.5em]
\texttt{cls.seq\_relationship.* | UNEXPECTED} &
Layer Next Sentence Prediction non usato. Scartato. \\[0.5em]
\texttt{embeddings.position\_ids | UNEXPECTED} &
Buffer deprecato nelle versioni recenti di Transformers. Ignorato. \\[0.5em]
\texttt{Unauthenticated requests to HF Hub} &
Modelli scaricati senza token HuggingFace. Funziona, ma con
rate limiting più restrittivo. Risolvibile impostando
\texttt{HF\_TOKEN}. \\
\bottomrule
\end{tabular}
\caption{Warning noti e loro interpretazione}
\end{table}

% 
\chapter{Glossario}
% 

\begin{description}[style=nextline]
  \item[A-Box (Assertional Box)]
    La parte del grafo RDF contenente gli individui (istanze), in
    contrapposizione alla T-Box che contiene le definizioni di classi
    e proprietà.

  \item[Anisotropy]
    Fenomeno per cui gli embedding prodotti da BERT raw tendono a
    occupare una regione ristretta dello spazio vettoriale, con tutte
    le direzioni simili tra loro. Causa score di cosine similarity
    sistematicamente alti per qualsiasi coppia di testi.

  \item[CLS token]
    Token speciale \texttt{[CLS]} inserito da BERT all'inizio di ogni
    sequenza di input. Il suo vettore nell'ultimo hidden state viene
    usato come embedding dell'intera sequenza.

  \item[Cosine Similarity]
    Misura di similarità tra due vettori basata sull'angolo che
    formano, indipendentemente dalla loro magnitudine.

  \item[Embedding]
    Rappresentazione vettoriale densa di un testo in uno spazio
    numerico ad alta dimensionalità, prodotta da un modello
    neuronale. Testi semanticamente simili hanno embedding vicini.

  \item[Entity Linking]
    Task di NLP che consiste nel collegare menzioni testuali di
    entità a voci corrispondenti in una base di conoscenza strutturata.

  \item[Fine-tuning]
    Processo di ri-addestramento di un modello pre-addestrato su un
    dataset specifico per un task specifico, aggiornando parzialmente
    i pesi del modello.

  \item[Knowledge Graph]
    Grafo strutturato che rappresenta entità e le relazioni tra esse
    in forma di triple soggetto-predicato-oggetto.

  \item[O*NET]
    Occupational Information Network: database del Dipartimento del
    Lavoro degli USA che fornisce descrizioni standardizzate di
    migliaia di professioni, incluse competenze, attività e titoli
    alternativi.

  \item[RDB2RDF]
    Standard W3C per la mappatura di dati relazionali (database SQL)
    in dati RDF. Include il Direct Mapping e R2RML.

  \item[T-Box (Terminological Box)]
    La parte del grafo RDF (o dell'ontologia OWL) contenente le
    definizioni di classi, proprietà e assiomi -- in contrasto con
    l'A-Box che contiene gli individui.

  \item[Triple RDF]
    Unità atomica del modello RDF, nella forma
    \textit{soggetto -- predicato -- oggetto}.
\end{description}

\end{document}